\documentclass[11pt,letterpaper]{article}
\usepackage{amsmath,amssymb,amsthm}
\renewcommand{\thefootnote}{\fnsymbol{footnote}}
\begin{document}
\title{Random Walks on the Integers\thanks{Notes prepared for a reading group.}}
\author{Ingrid Solheim \and Kofi Mensah}
\date{}
\maketitle
\begin{abstract}
We compute the probability that a simple random walk started at $k$ reaches $N$ before $0$, and the expected duration of the game.
\end{abstract}

\section{Gambler's ruin}
Let $S_0=k$ and $S_{n+1}=S_n\pm1$ with probabilities $p$ and $q=1-p$. Write $h(k)$ for the probability of hitting $N$ before $0$. First-step analysis gives
\begin{equation}
h(k)=p\,h(k+1)+q\,h(k-1),\qquad h(0)=0,\ h(N)=1.
\end{equation}
For $p=\tfrac12$ the solution is linear, $h(k)=k/N$.\footnote{This is also a consequence of the martingale property.}

For $p\ne\tfrac12$,
\[
h(k)=\frac{1-(q/p)^k}{1-(q/p)^N}.
\]

\section{Expected duration}
\begin{flushleft}
In the fair case the expected number of steps is $k(N-k)$.\footnote{Compare with the variance of Brownian motion.}
\end{flushleft}
\begin{flushright}
\emph{Quite surprisingly, it is quadratic in $N$.}
\end{flushright}
\end{document}
